\section{Endpoint simplicity and the sharp order of finite detection}
\label{sec:analytic-multiplicity-detection-equivalence}

The dominant-double hypothesis in NN-A027 is needed for its displayed
\emph{exact constants}, but not for a matching order of growth when
some endpoint is multiple. We prove this using an actual compact
unitary orbit and an explicit Schur isolation argument.

\begin{lemma}[A negative Hermitian orbital Fourier value]
\label{lem:analytic-negative-orbital}
Let $D\ne0$ be a Hermitian $m$ by $m$ matrix with trace zero and
spectrum invariant under multiplication by $-1$. Then
\[
 \phi_D(X)=\int_{U(m)}e^{i\operatorname{Tr}(DUXU^*)}\,dU,
                           \qquad X=X^*,
\]
is real and takes a strictly negative value. One may choose a
negative value at a diagonal matrix with distinct positive eigenvalues.
\end{lemma}

\begin{proof}
The orbit of $D$ equals that of $-D$. Haar invariance therefore makes
$\phi_D$ real. It is continuous, conjugation-invariant, and
$\phi_D(0)=1$. Give the real vector space $\mathcal H_m$ of Hermitian
matrices the Euclidean inner product $\operatorname{Tr}(XY)$ and
Lebesgue measure for an orthonormal basis; its dimension is $d=m^2$.
Fubini and the elementary Gaussian Fourier integral give, for every
$t>0$,
\begin{equation}\label{eq:analytic-orbital-Gaussian}
 \int_{\mathcal H_m}e^{-t\|X\|^2/2}\phi_D(X)\,dX
       =\left(\frac{2\pi}{t}\right)^{d/2}
                         e^{-\|D\|^2/(2t)}.
\end{equation}
The right side tends to zero as $t\downarrow0$, because $D\ne0$.
If $\phi_D$ were nonnegative everywhere, a fixed ball about zero on
which $\phi_D\ge1/2$ would give a positive lower limit for the left
side, a contradiction. Thus some $X$ has $\phi_D(X)<0$.
Diagonalize it, perturb its eigenvalues to be distinct while preserving
strict negativity, and add a scalar matrix large enough to make every
eigenvalue positive. This last operation does not change $\phi_D$
because $\operatorname{Tr}D=0$.
\end{proof}

\paragraph{Exact compact integral locator.}
We use the classical Harish--Chandra--Itzykson--Zuber formula with Haar
probability measure. If $A=\operatorname{diag}(a_1,\ldots,a_m)$ and
$B=\operatorname{diag}(b_1,\ldots,b_m)$ have distinct eigenvalues, then
\begin{equation}\label{eq:analytic-HCIZ}
 \mathcal I(A,B):=\int_{U(m)}e^{\operatorname{Tr}(AUBU^*)}\,dU
       =\left(\prod_{j=1}^{m-1}j!\right)
                    \frac{\det[e^{a_i b_j}]}{\Delta(a)\Delta(b)}.
\end{equation}
It extends to complex diagonal $A$ and coincident eigenvalues by
holomorphic continuation; the integral is entire in the matrix entries
on compact sets. A precise primary locator is J.-B.~Zuber,
\emph{Horn's problem and Harish-Chandra's integrals. Probability
density functions}, Ann. Inst. Henri Poincar\'e Comb. Phys. Interact.~5
(2018), 309--338, orbital definition (6), p.~311, formula (11), p.~312,
and constant (12), p.~313, DOI~10.4171/AIHPD/56,
\url{https://www.numdam.org/item/10.4171/aihpd/56.pdf}.
Only this exact classical identity is imported; the sign and isolation
arguments needed below are proved here.

\begin{lemma}[A whole coalescing cluster gives a bounded scaled-width witness]
\label{lem:analytic-cluster-Schur-witness}
Let an $\ell^1$ alphabet, closed under conjugation, have the following
form as $\varepsilon\downarrow0$:
\[
 b_1(\varepsilon)>\cdots>b_p(\varepsilon)>0,\quad
 \xi_j(\varepsilon)=b(\varepsilon)
                  e^{\varepsilon a_j(\varepsilon)}\quad(1\le j\le m),
                       \quad (\eta_j(\varepsilon))_{j\ge1}.
\]
Suppose $b>0$, $b\to b^0>0$,
$b_1^0>\cdots>b_p^0>b^0$, and
\[
 a_j(\varepsilon)=i d_j+O(\varepsilon),\qquad
 \sum_j a_j(\varepsilon)=0,
\]
where the $d_j$ are distinct real numbers, not all zero, with a
sign-symmetric multiset. Suppose the tail has uniformly bounded
$\ell^1$ norm, $\sup_j|\eta_j|\le c_0<b^0$, and its product
$C_\varepsilon(z)=\prod_j(1-\eta_jz)^{-1}$ converges locally uniformly
on $|z|<1/c_0$ to a product of positive real endpoint parameters.
The assertions about $p$ are vacuous when $p=0$.

Then there is a constant $C>0$ and a family of partitions
$\lambda(\varepsilon)$ of fixed length at most $q=p+m$, with
$\lambda_1\le C/\varepsilon+1$, such that
\[
                        s_{\lambda(\varepsilon)}<0
\]
for all sufficiently small $\varepsilon$. This is a negative
Toeplitz minor of order at most $C/\varepsilon+1$ in an index
window of size at most $C/\varepsilon+q$.
\end{lemma}

\begin{proof}
Put $D=\operatorname{diag}(d_1,\ldots,d_m)$. The preceding lemma
provides distinct numbers
$L_1>\cdots>L_m>0$ with
$\phi_D(\operatorname{diag}(L))<0$. Choose $C>L_1$ and $\eta>0$
so small that $L_m>2\eta$ and $C-L_1>2\eta$. Put
\[
 \mu_i=\lfloor L_i/\varepsilon\rfloor,\qquad
 n=\lceil C/\varepsilon\rceil,\qquad
 \lambda=(\underbrace{n,\ldots,n}_{p\text{ times}},
                                  \mu_1,\ldots,\mu_m).
\]
These are partitions for small $\varepsilon$, with
$\mu_m\ge\eta/\varepsilon$ and
$n-\mu_1\ge\eta/\varepsilon$. All row lengths are $O(1/\varepsilon)$.

First consider just the cluster, divided by its positive common scale
$b$. Write $A_\varepsilon=\operatorname{diag}(a_j(\varepsilon))$,
$\rho=(m-1,m-2,\ldots,0)$,
$U_\varepsilon=\operatorname{diag}(\varepsilon(\mu+\rho))$ and
$V_\varepsilon=\operatorname{diag}(\varepsilon\rho)$.
The Weyl alternant and \eqref{eq:analytic-HCIZ} give the \emph{exact}
identity
\begin{equation}\label{eq:analytic-Schur-orbital-limit}
 \frac{s_\mu(e^{\varepsilon a_1},\ldots,e^{\varepsilon a_m})}
      {s_\mu(1,\ldots,1)}
   =\frac{\mathcal I(A_\varepsilon,U_\varepsilon)}
          {\mathcal I(A_\varepsilon,V_\varepsilon)}
      \longrightarrow\phi_D(\operatorname{diag}(L))<0.
\end{equation}
Indeed the Vandermonde quotient
$\Delta(\mu+\rho)/\Delta(\rho)$ is the positive Weyl dimension,
and the remaining determinants give the displayed ratio. Its denominator
tends to one. Thus the cluster Schur value is negative, of magnitude
at least $c\,b^{|\mu|}s_\mu(1^m)$ for a fixed $c>0$.
This conclusion is stable under the full $O(\varepsilon)$ perturbation
of $A_\varepsilon$; no exact geometric spacing of the cluster is used.

We now show that the finite positive prefix and infinite tail change
this sign only by a positive factor and an exponentially small error.
All constants below may depend on the fixed $p,m,L,C$ and the stated
spectral separations. Choose $c_1$ with $c_0<c_1<b^0$.
The complete-symmetric generating function has exactly $q=p+m$
simple top poles $\alpha\in\mathcal A=\{b_i,\xi_j\}$ inside
$|z|=1/c_1$. The cluster pole differences are bounded below by a
constant times $\varepsilon$, since the $d_j$ are distinct. The
partial-fraction argument of NN-A027 therefore gives
\begin{equation}\label{eq:analytic-full-cluster-poles}
 h_r=\sum_{\alpha\in\mathcal A}A_\alpha\alpha^r+g_r,
 \qquad |A_\alpha|\le C_2\varepsilon^{-(m-1)},\qquad
 |g_r|\le C_2\varepsilon^{-(m-1)}c_1^r.
\end{equation}
The analytic remainder bound follows on the same fixed circle after
removing all top poles. Uniform tail product bounds ensure its constant
is independent of $r$ and $\varepsilon$.

Apply \eqref{eq:analytic-full-cluster-poles} to the ordinary
Jacobi--Trudi determinant $\det[h_{\lambda_i-i+j}]_{i,j=1}^q$.
The part using all $q$ top poles is exactly
\begin{equation}\label{eq:analytic-full-cluster-leading}
                  s_\lambda(\mathcal A)
                              \prod_{\alpha\in\mathcal A}C_\varepsilon(1/\alpha).
\end{equation}
This is the same internal Vandermonde cancellation as in NN-A027,
now with a general partition instead of a rectangle. The product on
the right is real and tends to a positive number.

Here is a uniform error estimate for all terms with remainder columns.
Let $P=\prod_{i=1}^p b_i$ and $W=P^n b^{|\mu|}>0$.
A term with $r\ge1$ remainder columns can use only $q-r$ distinct
top pole labels; repeated labels give zero. After determinant expansion,
its exponential weight assigns the corresponding moduli, and $r$ copies
of $c_1$, to the nonincreasing row exponents $\lambda_i$.
The rearrangement inequality maximizes that weight by matching their
orders. Compared with the full set of top moduli, replacing any one
by $c_1$ loses a factor at most
$\widetilde q^{\mu_m}$ for a fixed $\widetilde q<1$:
all exponents are at least $\mu_m$, and every top modulus is uniformly
larger than $c_1$. More explicitly the sorted logarithms decrease
coordinatewise, their total decrease is at least
$\log(\min|\alpha|/c_1)$, and multiplication by exponents at least
$\mu_m$ gives this loss. The full top weight is bounded by a constant
times $W$, because $|\xi_j|/b=e^{O(\varepsilon^2)}$ and
$|\lambda|=O(1/\varepsilon)$. Fixed shifts $j-i$ in the matrix entries
affect only a constant. There are finitely many terms, and the residue
bounds in \eqref{eq:analytic-full-cluster-poles} introduce only a fixed
power of $\varepsilon^{-1}$. Thus the total remainder is
\begin{equation}\label{eq:analytic-tail-exponential-error}
                  O\bigl(W\varepsilon^{-K}
                              \widetilde q^{\mu_m}\bigr)
\end{equation}
for some fixed $K$. This explicitly avoids a determinant estimate that
would repeat the largest pole in all rows.

To isolate the prefix inside $s_\lambda(\mathcal A)$, expand its Weyl
numerator along the first $p$ rows. The leading choice assigns the
$p$ positive prefix columns to those rows. Dividing by the full
Vandermonde gives exactly
\begin{equation}\label{eq:analytic-prefix-factorization}
 P^n s_\mu(\xi_1,\ldots,\xi_m)
            \prod_{i=1}^p\prod_{j=1}^m(1-\xi_j/b_i)^{-1}.
\end{equation}
For example the prefix numerator block is
$P^{n+m}\Delta(b_1,\ldots,b_p)$, and the cluster block is
$s_\mu(\xi)\Delta(\xi)$; these cancel the corresponding full
Vandermonde factors. The remaining product is real and tends to a
positive number.

Every other prefix assignment puts at least one larger $b_i$ into a
cluster row, whose length is at most $\mu_1$, and a cluster modulus
into a row of length $n$. It therefore loses a factor at most
$q_2^{n-\mu_1}$, with $q_2<1$ fixed. Pair the misplaced prefix and
cluster columns; the ratio $|\xi|/b_i$ is uniformly less than one,
and apply the rearrangement inequality to any other columns.
The cluster Vandermonde in the denominator has magnitude bounded
below by a constant times $\varepsilon^{m(m-1)/2}$, while all other
Vandermonde factors stay away from zero. Consequently the total wrong
assignment contribution is
\begin{equation}\label{eq:analytic-prefix-exponential-error}
                   O\bigl(W\varepsilon^{-K'}q_2^{n-\mu_1}\bigr).
\end{equation}
For $p=0$ this error is absent. Combining
\eqref{eq:analytic-full-cluster-leading}--\eqref{eq:analytic-prefix-exponential-error},
the full Schur value is a positive factor, bounded away from zero,
times $P^n s_\mu(\xi)$, plus $W$ times an exponentially small quantity
in $1/\varepsilon$. Equation~\eqref{eq:analytic-Schur-orbital-limit}
therefore makes it negative eventually.

Finally put $r=\lambda_1\le n$; equality need not hold when $p=0$.
Dual Jacobi--Trudi gives a minor of order $r$:
\[
                     s_\lambda=\det[e_{\lambda'_i-i+j}]_{i,j=1}^r.
\]
Transpose this determinant. Use row indices $r-i$ and column indices
$\lambda'_j+r-j$, both in decreasing order, and then reverse both.
The resulting increasing index sets lie in $0,\ldots,r+q-1$,
which is contained in $0,\ldots,n+q-1$.
This proves the stated order and window bounds.
\end{proof}

\begin{theorem}[NN-A028: a finite-detection characterization of endpoint simplicity]
\label{thm:analytic-simplicity-detection-equivalence}
Assume A. If $H_0$ has any multiple zero, then there are positive
constants $c,C,a_0$ such that for $0<a<a_0$,
\begin{equation}\label{eq:analytic-any-multiple-detection-scale}
 c a^{-1/2}\le m_*(-a),\ R_*(-a),\ d_*(-a)\le C a^{-1/2}.
\end{equation}
Here the three quantities are respectively the least negative-minor
order, least detecting index window for $F_{-a}$, and least failing
rank of the actual translated unitary characters from NN-A021.
The chained notation means that \emph{each} of the three quantities
lies between the displayed bounds.

Consequently the following four statements are equivalent:
\begin{enumerate}
\item Every nontrivial zero of $\zeta$ is simple.
\item $\sqrt a\,m_*(-a)\longrightarrow+\infty$ as $a\downarrow0$.
\item $\sqrt a\,R_*(-a)\longrightarrow+\infty$ as $a\downarrow0$.
\item $\sqrt a\,d_*(-a)\longrightarrow+\infty$ as $a\downarrow0$.
\end{enumerate}
The equivalence proves none of these equivalent assertions.
When a multiple endpoint exists, the upper witnesses have partition
height bounded by a fixed number depending on its first occurrence;
only their width and corresponding unitary rank must grow.
NN-A027 supplies sharper exact constants in its dominant-double case.
\end{theorem}

\begin{proof}
Suppose a multiple positive endpoint exists and choose the smallest
one, $x_0$, of multiplicity $m\ge2$. Such a first one exists because
the endpoint divisor is discrete and bounded away from zero. All
$p$ positive endpoint zeros below it are simple, so their inverse-square
parameters are a strictly separated positive prefix for sufficiently
small negative time. The local Hermite expansion from NN-A006 gives
\[
             z_j(-a)=x_0+i\sqrt{2a}\,h_{m,j}+O(a),
                             \qquad 1\le j\le m.
\]
These local roots are distinct for $a>0$ small. Their inverse squares
form a conjugation-stable cluster. Take its positive geometric mean
$b(a)=(\prod_j z_j(-a)^{-2})^{1/m}$ and put $\varepsilon=\sqrt a$.
Branches of the logarithm near the positive endpoint give
\[
 z_j(-a)^{-2}=b(a)e^{\varepsilon a_j(\varepsilon)},\qquad
 a_j(\varepsilon)=-\frac{2i\sqrt2}{x_0}h_{m,j}+O(\varepsilon),
                         \qquad\sum_j a_j(\varepsilon)=0.
\]
The limiting Hermite diagonal is nonzero, traceless, sign-symmetric
and has distinct entries. The uniform zero-count and tail argument
in NN-A027 gives the required smaller-modulus infinite tail and its
locally uniform positive endpoint product. Therefore
Lemma~\ref{lem:analytic-cluster-Schur-witness} supplies a negative
minor of order at most $C/\sqrt a+1$ and window at most
$C/\sqrt a+p+m$. This proves the upper bounds for $m_*$ and $R_*$.

For the actual unitary character use, as in NN-A027,
\[
       \widetilde\beta_j(a)=\frac1{p_a z_j(-a)^2-1},
                             \qquad p_a=F_{-a}'(0)>0.
\]
The positive endpoint ordering is unchanged. The cluster has the
same Hermite diagonal multiplied by the positive factor
$(1-1/(p_0x_0^2))^{-1}$. Its tail is uniformly $\ell^1$ and separated,
since the normalized original parameters have modulus bounded by
some fixed number less than one near $a=0$. The same lemma thus
gives a negative Toeplitz minor of order at most $C'/\sqrt a+1$
for the translated coefficient sequence. The exact finite-rank
criterion in NN-A021 turns it into failure of positive type at that
rank. The dual partition specifies the actual negative irreducible
coefficient; its entries are bounded by the fixed height $p+m$.

The positive lower bounds of inverse-square-root order are already
proved in NN-A022, including its explicit order/window bound and
its effective unitary-rank bound. Decreasing $c$ and $a_0$ absorbs
the fixed additive constants, proving
\eqref{eq:analytic-any-multiple-detection-scale}.

If every endpoint is simple, NN-A022 instead proves
$\sqrt a\,m_*(-a)\to\infty$ and
$\sqrt a\,d_*(-a)\to\infty$; the elementary inequality
$R_*\ge m_*-1$ gives the window statement. Conversely, the existence
of a multiple endpoint gives the bounded upper quantities just proved,
so none of these three limits can be infinite. A identifies the
positive endpoint zeros with the positive nontrivial zeta zeros,
including their multiplicities. This proves the four-way equivalence.
Rodgers--Tao and A ensure the least witnesses are finite at every
negative time even in the all-simple case; they do not decide which
case occurs.
\end{proof}

\paragraph{Geometric and representation meaning.}
The first nonreduced endpoint local algebra has finite length $m$.
Its universal leading heat unfolding is the Hermite cluster, whose
nonzero traceless Hermitian diagonal has the genuine compact orbit
$U(m)/U(1)^m$. This orbit is a full flag manifold with its usual
coadjoint-orbit symplectic form. The Fourier transform of its Haar
orbital measure is forced to become negative by
\eqref{eq:analytic-orbital-Gaussian}. The exact Weyl--HCIZ identity
then produces negative values of actual integral highest-weight
Schur characters. The finite positive prefix and infinite analytic
tail have been isolated with explicit uniform exponentially small
errors, so the same sign occurs in the actual heat coefficient minor.
The exterior-power and representation-ring operations have the exact
locator SGA~6, Expos\'e~V, Example~2.2(b) and Lemma~2.2.1.
No integrality or prequantization of the Hermite orbit's real symplectic
class is asserted. The result is an exact relation between reducedness
of the endpoint analytic divisor and growth of three concrete finite
certificate families; it is not a proof of reducedness, an arithmetic
purity realization, or a complexity-class separation.
